% Generated by scripts/tikz_reviews.py from the figure manifests.
\ReviewFigure{Interior points and failure of exact recovery}
{figures/sparse-theory/proofs/polytope-proof-1.png}
{figures/tikz/sparse-theory-polytope-proof.pdf}
{The onto-map polytope characterization: x0 solves basis pursuit iff Ax0/||x0||1 lies on the boundary of AB.}
{Uses the normalized polytope AB and q=Ax0/r, r=||x0||1, consistently with the corrected proof. Replaces the original ambiguous label z on an observation-space boundary point by (1+t)q. Both panels illustrate failure of optimality through interior membership; no uniqueness claim is implied.}
{sparse-theory-polytope-proof}
\ReviewFigure{Removing a dependent active coefficient}
{figures/sparse-theory/proofs/injectivity.png}
{figures/tikz/sparse-theory-injectivity.pdf}
{Support reduction along x(t)=x+t h when A\_I h\_I=0.}
{The original draws affine coefficient trajectories. Three representative coordinates are drawn with a marked sign-preserving interval. At its left endpoint the first coefficient vanishes. The constant-fidelity and affine-l1 statements follow the chapter proof; the drawn slopes also sum to zero within the positive orthant.}
{sparse-theory-injectivity}
\ReviewFigure{Bregman divergence as a vertical gap}
{figures/sparse-theory/proofs/bregman.png}
{figures/tikz/sparse-theory-bregman.pdf}
{Definition of D\_eta(x|x0) for eta in the subdifferential of a convex J.}
{The left panel uses a strictly convex quadratic and its tangent; the positive vertical gap is the Bregman divergence. The right panel uses J=abs(x), with x0 and x on the same positive affine branch, so the divergence vanishes even when x differs from x0.}
{sparse-theory-bregman}
\ReviewFigure{A strict subgradient margin controls the error}
{figures/sparse-theory/proofs/bregman-l1.png}
{figures/tikz/sparse-theory-bregman-l1.pdf}
{The coordinatewise estimate in the l1 Bregman lower bound.}
{Explicitly sets x0=0, as required on nonsaturated coordinates for J=abs(x). For abs(eta)<1 the supporting line eta*x lies strictly below abs(x) away from zero. The lower bound is D\_eta(x|0)>= (1-abs(eta))*abs(x), not a norm bound on saturated coordinates.}
{sparse-theory-bregman-l1}
\ReviewFigure{Parameter region for sign consistency}
{figures/sparse-theory/proofs/recov-zone.png}
{figures/tikz/sparse-theory-recovery-region.pdf}
{Intersection of delta+lambda<T/K and R*delta<S*lambda in the support-stability proof.}
{Reconstructs the two strict inequalities and the contained simpler triangle. The vertical cutoff is lambda\_max=TR/[K(R+S)], using R=KL+1. Shading indicates interiors; boundary lines are not included in the sufficient conditions. The axes follow the caption: horizontal lambda, vertical delta.}
{sparse-theory-recovery-region}
\ReviewFigure{Convolution of a discrete measure}
{figures/sparse-theory/proofs/spikes.png}
{figures/tikz/sparse-theory-convolution-spikes.pdf}
{The dictionary map phi(z), discrete measure m\_x, and measurement operator A.}
{Two shifted kernels and two point masses are illustrative; their locations are reconstruction-grid points z\_j and z\_i. The output is their weighted sum, written as the continuous operator applied to m\_x and as the discretized dictionary product Ax. The diagram distinguishes the index-space weights x from the physical locations z.}
{sparse-theory-convolution-spikes}
